\chapter{Conclusion}\label{chap:conclusion}

The preliminary results presented in this qualification demonstrate that grounding image analysis in the physical characteristics of the observed phenomenon, rather than relying on brute-force feature learning, leads to a detection pipeline that is both accurate and transparent. The preprocessing stage effectively isolates the plasma depletion signatures from environmental contamination, and the compact feature representation produced by 2DPCA retains enough structural information for the ensemble classifier to distinguish between EPB events and clear-sky conditions with high confidence. The interpretability analysis further reinforced this outcome by confirming that the classifier's attention aligns with the expected spatial distribution of ionospheric depletions.

These findings establish a solid foundation for the next phases of this dissertation. Moving forward, a central priority will be assessing how well this architecture generalizes to observations from different geographical stations and across varying solar cycle conditions, where EPB morphology and occurrence rates differ considerably. We also plan to investigate the integration of complementary data sources, such as GNSS scintillation indices and radar backscatter measurements, to provide a more comprehensive view of ionospheric irregularities. Finally, exploring alternative and more recent machine learning architectures will allow us to benchmark the current approach and identify potential improvements for operational deployment.
